% SPDX-License-Identifier: Apache-2.0
%
% HDMI in TxHDL. Built by //docs:hdmi. Every listing is included from
% the tree, and every printout and figure is what the build made by
% running the example.
\documentclass[journal,9pt]{IEEEtran}

\newcommand{\listingsize}{\fontsize{6}{7}\selectfont}
% SPDX-License-Identifier: Apache-2.0
%
% The house preamble, shared by every document in this directory. Loaded
% after \documentclass, before \title. Keeping it in one file is what
% stops two articles drifting apart in font, listing style or figure
% style.
% Computer Modern. IEEEtran selects Times (ptm) by default, so the face has
% to be chosen explicitly.
%
% Latin Modern is the Computer Modern variety used here, and the choice is
% forced rather than aesthetic. Under T1 encoding, plain `cmr` has no Type 1
% outlines unless cm-super is installed, and it is not in the pinned
% distribution, so pdflatex falls back to EC bitmaps and every glyph in the
% PDF becomes a Type 3 font. That was measured: the first build produced a
% document with 10 Type 3 fonts and no Type 1. Latin Modern is Computer
% Modern redrawn with a full T1 repertoire and real .pfb outlines, so the
% same design comes out scalable.
\usepackage[T1]{fontenc}
\usepackage[utf8]{inputenc}
\usepackage{lmodern}
% microtype is not in the pinned TeX distribution. emergencystretch alone
% keeps the two-column measure from producing overfull lines.
\setlength{\emergencystretch}{3em}

\usepackage{amsmath}
\usepackage{amssymb}
\usepackage{amsthm}
\usepackage{booktabs}
\usepackage{array}
% enumitem is not in the pinned TeX distribution; the list spacing below
% does what \setlist would have done.
\usepackage{float}
\usepackage{xcolor}
\usepackage{listings}
\usepackage{tikz}
\usetikzlibrary{arrows.meta,positioning,fit,backgrounds,calc}
\usepackage[hidelinks,bookmarks=true]{hyperref}

% Numbered, definition-style box for a rule the reader is meant to apply.
\theoremstyle{definition}
\newtheorem{rulebox}{Rule}
\newtheorem{example}{Example}

% Unnumbered asides.
\theoremstyle{remark}
\newtheorem*{tip}{Tip}
\newtheorem*{pitfall}{Pitfall}

% Tighter lists, in place of enumitem's \setlist.
\setlength{\topsep}{2pt}
\setlength{\itemsep}{1pt}
\setlength{\parsep}{0pt}

\newcommand{\code}[1]{\texttt{#1}}
\newcommand{\term}[1]{\emph{#1}}

% Syntax highlighting. Four roles, distinguishable in colour and still
% distinguishable in greyscale, because an IEEE article gets printed: the
% keyword is the only bold role, the comment is the only italic one, and the
% two remaining roles differ in lightness as well as in hue.
\definecolor{kwblue}{RGB}{20,60,150}
\definecolor{cmtgreen}{RGB}{90,120,90}
\definecolor{strred}{RGB}{150,50,40}
\definecolor{numgray}{gray}{0.45}
\definecolor{framegray}{gray}{0.70}
\definecolor{bgshade}{gray}{0.97}
% A darker fill for figure cells. The listing background has to stay very
% light to keep code readable; a figure cell has to be clearly shaded or the
% distinction it draws is invisible in print.
\definecolor{cellshade}{gray}{0.90}

% The listing size is the document's choice, because it depends on the
% column: a two-column article sets `\listingsize` to 6pt and keeps its
% listings within 70 columns, a one-column document keeps the body size
% and includes source files at 80. Lines are never broken; a line that does not fit
% overflows the frame, which is visible, rather than wrapping, which is
% not.
\providecommand{\listingsize}{\scriptsize}

% `'` is deliberately not a string delimiter: the language uses it for loop
% labels, as Rust does, so treating it as a quote would swallow the rest of
% the line.
\lstdefinestyle{txhdl}{
  basicstyle=\ttfamily\listingsize,
  keywordstyle=\bfseries\color{kwblue},
  commentstyle=\itshape\color{cmtgreen},
  stringstyle=\color{strred},
  numberstyle=\tiny\color{numgray},
  morekeywords={mod,use,pub,const,type,struct,enum,trait,impl,fn,async,let,mut,
    match,if,else,loop,while,for,in,break,continue,return,as,await,
    module,bus,channel,transaction,protocol,pipeline,flow,seq,config,
    port,stage,pipe,instance,connect,bind,tag,domain,blackbox,foreign,
    property,role,signal,handler,true,false,
    sequence,interface,modport,implements,package,import,var,wire,func,proc,
    case,default,next,fork,branch,join,state,fsm},
  morecomment=[l]{//},
  morestring=[b]",
  showstringspaces=false,
  columns=fullflexible,
  keepspaces=true,
  breaklines=false,
  % Framed, so a listing is bounded rather than floating in the column.
  frame=single,
  rulecolor=\color{framegray},
  framesep=3pt,
  backgroundcolor=\color{bgshade},
  xleftmargin=3pt,
  xrightmargin=3pt,
  aboveskip=6pt,
  belowskip=6pt,
  % Every listing is numbered and captioned, so it can be referred to by
  % number. `caption` without `float` numbers the listing and leaves it
  % where it was written, which is what a listing in running prose needs.
  captionpos=b,
  abovecaptionskip=2pt,
  belowcaptionskip=6pt,
}

% Figures drawn as text. Framed the same way, and with the highlighting
% switched off, because a box-and-arrow diagram is not code and half its
% words turning blue reads as an accident. Setting a style adds keys rather
% than clearing the ones already set, so the three roles are neutralised by
% hand rather than by leaving them out.
\lstdefinestyle{ascii}{
  basicstyle=\ttfamily\listingsize,
  keywordstyle=\ttfamily,
  commentstyle=\ttfamily,
  stringstyle=\ttfamily,
  columns=fullflexible,
  keepspaces=true,
  frame=single,
  rulecolor=\color{framegray},
  framesep=3pt,
  backgroundcolor=\color{bgshade},
  xleftmargin=3pt,
  xrightmargin=3pt,
  aboveskip=6pt,
  belowskip=6pt,
}
\lstset{style=txhdl}

% House TikZ styles. `shorten` lives on the arrow styles rather than on each
% arrow, so every arrow in the document gets the gap, including ones added
% later. TikZ clips a path at a node's border, so without it an arrowhead
% butts against the box with no gap at all. Keep `node distance` well above
% twice the shorten, or the arrow shrinks to nothing.
\tikzset{
  box/.style={draw,rounded corners=2pt,align=center,inner sep=4pt,
              font=\footnotesize},
  boxf/.style={box,fill=bgshade},
  lbl/.style={font=\scriptsize\itshape,align=center},
  fl/.style={-{Stealth[length=4pt]},shorten >=2.5pt,shorten <=2.5pt},
  fld/.style={fl,dashed},
}



\InputIfFileExists{buildstamp.tex}{}{}
\providecommand{\buildstamp}{Draft build (outside the Bazel flow).}

% A range of the HDMI part, by its begin/end markers.
\newcommand{\hdmipart}[3]{%
\lstinputlisting[rangebeginprefix=//\ begin\{,rangebeginsuffix=\},%
rangeendprefix=//\ end\{,rangeendsuffix=\},includerangemarker=false,%
linerange=#1-#1,caption={#2},label={#3}]{lib/parts/src/hdmi.rs}}

\title{HDMI in TxHDL}

\author{Filip Filmar and Dragiša Janković%
\thanks{Both authors are independent. Filip Filmar is the
corresponding author, at f@hdlfactory.com; Dragiša Janković is
reached at linkedin.com/in/dragisa-jankovic.}%
\thanks{This is the tutorial on the HDMI part, for a reader who knows
the AXI link and its AXI-Lite bridge~\cite{axi}. Every listing is
included from the project repository \code{HDL/txhdl} at build time,
and every printout and figure is what the build produced by running
the example. The cover~\cite{cover} lists every document.}%
\thanks{The authors used a large language model, Claude, as an
assistant in exploring the concepts and in writing and constructing
the documents and the programs; every commit of the repository says
so and carries the prompts that led to it, verbatim.}%
\thanks{\buildstamp}}

\markboth{HDMI in TxHDL}{Filmar and Janković: HDMI in TxHDL}

\begin{document}
\maketitle

\begin{abstract}
The Alinx AX7A200B sends HDMI through a SiI9134 encoder. The chip
does the TMDS encoding and the serialising. The FPGA gives it a pixel
clock, 24 bits of colour, two syncs and a data enable, and a few
register writes over I2C before it shows anything. This part is
two units in the lowered subset. The video peripheral counts the
raster, shows a framebuffer of 12-bit pixels, and lets a host paint
that framebuffer over AXI-Lite. The I2C master resets the chip and
writes its configuration. The build runs both, in a mode of 16 by 12
pixels, against a host that paints a picture and a model of an I2C
device, checks every pixel on the screen, and simulates both netlists
against the run under \code{nvc} and under Verilator. For the board
there is a hand-written demonstration that shows a test picture at
640 by 480 and 60~Hz, through Vivado to a bitstream; it has not yet
run on the board.
\end{abstract}

\section{What Is Here}
\label{sec:h-what}

The part is \code{hdmi} in \code{txhdl\_parts}, and it has two units.

\code{Hdmi} is the video peripheral. It runs on the pixel clock and
drives the chip's inputs: \code{rgb}, 24 bits with red on top,
\code{hsync} and \code{vsync}, low while active, and \code{de}, high
over the visible part of the raster. Its framebuffer holds 12-bit
pixels, four bits each of red, green and blue, and each framebuffer
pixel covers a square of \code{1 << SHIFT} screen pixels on a side. A
host writes the framebuffer through the unit's AXI-Lite port.

\code{I2cInit} is the chip's setup. It holds the chip's reset low,
lets it go, waits, and writes a table of registers as I2C
transactions. It reports \code{done} when the table is written and
\code{failed} when the chip did not acknowledge a byte.

\section{The Raster}
\label{sec:h-raster}

A frame is a raster of columns and rows, most of them visible and the
rest blanking: a front porch, a sync pulse and a back porch, across
and down. At 640 by 480 and 60~Hz that is 800 columns by 525 rows, a
pixel every 39.7~ns at 25.175~MHz, with both syncs low while active.
Listing~\ref{lst:h-modes} states those numbers. The unit takes all
eight of them as type parameters, so a test can run a mode of a few
hundred pixels.

\hdmipart{modes}{640 by 480 at 60~Hz.}{lst:h-modes}

The unit counts the column in \code{hc} and the row in \code{vc}. From
them it works out whether the beam is in the visible part and whether
it is in either sync pulse, with the functions of
Listing~\ref{lst:h-fns}. It reads the framebuffer at the pixel under
the beam into \code{px} at the clock edge, the synchronous read a block
RAM has. The syncs and the enable go into registers at the same edge,
so they meet the pixel they belong to. The chip's inputs are those
registers.

Three of the functions serve both axes, each called once per axis,
because the lowering binds a function's constant parameters to the
arguments of the call by position~\cite{i127}; the names inside the
function are its own. Until that fix the binding was by name, so a
function's parameters had to be named as the unit named its own, and
the functions came in pairs where one would have done. The last
column and the last row still have a function each: they count in
eight bits and in seven.

\hdmipart{fns}{The raster's functions.}{lst:h-fns}

\section{The Framebuffer and the Bus}
\label{sec:h-bus}

The framebuffer is a memory of $2^{15}$ words, addressed by the row
in the top seven bits and the column in the low eight. At 640 by 480
with \code{SHIFT} of 2 the picture is 160 by 120, and the rest of the
memory is unused.

\begin{table}[htbp]
\caption{The video peripheral's words.}
\label{tab:h-words}
\centering\footnotesize
\begin{tabular}{@{}lll@{}}
\toprule
Word & Access & Bits \\
\midrule
0, status & read & 0 vertical blank, 31 to 16 frames \\
1, cursor & read, write & 7 to 0 column, 14 to 8 row \\
2, pixel & write & 11 to 0 the colour \\
\bottomrule
\end{tabular}
\end{table}

A write to the pixel word puts the colour into the framebuffer at the
cursor and moves the cursor to the next pixel. At the end of a row the
cursor goes to the start of the next row, and after the last row back
to the first. A host sets the cursor once and then paints a run of
pixels in one fixed burst to the pixel word, since the AXI-Lite bridge
makes each beat of a fixed burst a transaction of its own at the same
address.

\hdmipart{state}{The video peripheral's state.}{lst:h-state}

\hdmipart{run}{The video peripheral's step.}{lst:h-run}

The colour going out is built from wires named for each colour, with a
function call per wire before the concatenation. A chain that begins
with a call of a function is not read by the lowering, which takes it
for a constant~\cite{i128}.

\section{Configuring the Chip}
\label{sec:h-i2c}

I2C is two open-drain lines, a clock and a data line, each pulled high
on the board and pulled low by whoever is talking. A transaction starts
with the data line falling while the clock is high and stops with it
rising while the clock is high. Between, each byte is eight bits, most
significant first, each changing while the clock is low and read while
it is high, and after each byte the device pulls the data line low
through one more clock to acknowledge it.

The master cuts a bit into four quarters of \code{DIV} cycles each,
the clock low in the first and the last. A transaction is 29 steps of
four quarters: the start, the device's address and its acknowledge, the
register and its acknowledge, the value and its acknowledge, and the
stop. \code{scl\_low} and \code{sda\_low} high pull a line low, and the
data line comes back in on \code{sda\_in}, where the master reads it in
the third quarter of each acknowledge.

\hdmipart{i2cstate}{The master's state.}{lst:h-i2cstate}

\hdmipart{i2c}{The master's step.}{lst:h-i2c}

Listing~\ref{lst:h-table} is the table the master writes: a register
that takes the chip out of power-down for the input the board wires,
and one that sets its output to DVI, the picture with no HDMI packets.
None of the documents this part was written from states the chip's
registers, so these values are to be confirmed against the SiI9134's
documentation or the board vendor's example before the board is
expected to show a picture. They are one table in the source, and the
functions that give the master its entries.

\hdmipart{table}{The chip's configuration, to be
confirmed.}{lst:h-table}

\section{The Run}
\label{sec:h-run}

\code{ex\_hdmi} puts the video peripheral behind \code{LiteBridge1},
behind a host's tracker, in a mode of 16 by 12 visible pixels with
each framebuffer pixel 2 by 2, so the framebuffer is 8 by 6. The host
client sets the cursor to the corner and paints all 48 pixels in one
fixed burst, red rising across, green rising down, and blue on the
diagonal. Then it reads the status until two more frames have been
shown.

Beside it the I2C master runs against \code{I2cDevice}, a model of a
device on the lines that acknowledges every byte and records what it
was sent. The run checks that it was sent the table.

The run rebuilds each frame from the pixels \code{de} marks, starting
at the vertical sync, and checks every one of the 192 screen pixels of
the last whole frame against the colour painted. A pixel shifted by
one column or one row would fail at the edge of its 2 by 2 square.

\begin{figure*}[t]
\centering
\input{hdmi_timing.tex}
\caption{The run: the pixels written on \code{w}, the syncs and the
enable frame after frame, and the chip's reset and
the I2C lines while the two transactions go out, until \code{done}.}
\label{fig:h-run}
\end{figure*}

\lstinputlisting[caption={\code{ex\_hdmi.rs}. Run with
\code{bazel run //lib/examples:ex\_hdmi}.},%
label={lst:h-ex}]{lib/examples/ex_hdmi.rs}

\lstinputlisting[style=ascii,caption={What \code{ex\_hdmi} printed when
the build ran it: the two transactions the device was sent, and the
last whole frame, one screen pixel per framebuffer pixel.},%
label={lst:h-out}]{out_hdmi.txt}

The peripheral and the master are lowered, and the build simulates
each netlist against this run under \code{nvc} and under Verilator, at
its ports. Tests in the part check the syncs and the enable against
the mode, one fixed lag fitting all three, and the table on the I2C
lines.

\section{On the Board}
\label{sec:h-board}

The board's part is written by hand in \code{//hdmi}. It is a
demonstration with no host: the netlist the build writes for it has a
test picture in the framebuffer's initial values, eight colour bars
over a grey ramp inside a white border, and the peripheral's AXI-Lite
inputs are held idle.

\lstinputlisting[caption={\code{hdmi/src/netlist.rs}: the netlist, with
the test picture.},label={lst:h-netlist}]{hdmi/src/netlist.rs}

The pixel clock comes from an MMCM. The board's 200~MHz input is
divided by ten and multiplied by 63 for a 1260~MHz oscillator, which
divided by 50 is 25.2~MHz, within half a per cent of the mode's
25.175~MHz. The chip's clock is the pixel clock turned over by an
\code{ODDR}, so the chip's rising edge falls in the middle of each
pixel. The I2C lines are open drain: the pad is driven low while the
master pulls, and left floating otherwise.
The two units' reset is the clock manager's lock, low until the pixel
clock is good, through two flip-flops on that clock, so that the
release lands on an edge of the clock it releases. A lowered module's
reset arm runs from its \code{rst} input and from nothing else, so a
top that left the input unconnected had the units start from the
power-on values alone and never reset them (issue 509).

\lstinputlisting[caption={\code{hdmi/board/hdmi\_demo.v}: the
demonstration.},label={lst:h-demo}]{hdmi/board/hdmi_demo.v}

\code{bazel build //hdmi:demo\_synth} puts the two lowered units and
the demonstration through Vivado's synthesis for the board's part, with
no error and no critical warning. Its 87 warnings are all the same
one: the bits of the peripheral's AXI-Lite inputs that the
demonstration holds at zero have no load. \code{bazel build
//hdmi:demo\_pnr} places and routes the design with the board's pins
and writes the bitstream. Every timing constraint is met, with
35.872~ns of setup slack and 0.164~ns of hold slack over 398
endpoints, and the design rule check finds nothing. The design takes
60 LUTs, 86 flip-flops, 12 block RAM tiles for the framebuffer, and 36
pins. No input or output delay is constrained on the chip's pins, so
that slack says nothing about the timing between the FPGA and the
chip.

The constraints name the clock, the LEDs, and the chip's 31 pins from
the table in the AX7A200B manual.

\lstinputlisting[caption={\code{hdmi/board/ax7a200\_hdmi.xdc}: the
board's pins.},label={lst:h-xdc}]{hdmi/board/ax7a200_hdmi.xdc}

\section{What Is Not Done}
\label{sec:h-not}

The demonstration has run on the board. On September 19, 2026 a monitor
on the board's HDMI connector showed the test picture as the
framebuffer holds it: eight colour bars over the top two thirds, white,
yellow, cyan, green, magenta, red, blue and black, a grey ramp beneath
them, and a white border. The chip's LEDs read locked and configured,
with no byte refused, so every register write was acknowledged.

That answers the three assumptions this section used to list. The
register values are right, since the chip accepted them and put a
picture out. The order of the data inputs is right, since the bars came
out in the order the framebuffer holds them and a swapped red and blue
would have turned yellow into cyan. The I/O standard is LVCMOS33, and
the vendor's own demonstration for this board uses it on the same pins.

One thing the board did add. The vendor's demonstration drives two
reset outputs, on L18 and on Y17, the second named for revision 1.0 of
the board, and this design drove only Y17, which the manual's table
names. It now drives both, since which ball reaches the chip depends on
the revision and a wrong guess holds the chip in reset with nothing to
say why.

What is still not done on the board side: the timing at the pins is not
constrained, so nothing has been measured about the clock's edge against
the data's at the connector, and the picture arriving says only that the
margin is not zero.

The whole peripheral runs on the pixel clock, bus port included. A
host on the system clock reaches it through a crossing of the five
AXI-Lite channels, which the demonstration does not need and nothing
here provides. Such a crossing written in the language is
\code{ChanCdc}, the clock-crossing channel in \code{//lib/parts}
under \code{cdc}, and the lowering co-simulates a unit of two clocks
against one run, each port checked on the edge of its own clock, since
issue 131~\cite{i131}; the flagship crosses these five channels with
the hand-written \code{chan\_cdc} instead, five instances in its top.

The framebuffer is 12 bits a pixel and at most 256 by 128. The mode
is fixed when the type is named, and there is one I2C table.

\begin{thebibliography}{9}

\bibitem{axi}
F.~Filmar and D.~Janković, ``AXI in TxHDL,'' project repository
\code{HDL/txhdl}, \code{//docs:axi}, 2026.

\bibitem{i127}
Issue 127, ``fix(lower): a function's const parameters are bound by
name, not by position,'' project repository \code{HDL/txhdl}, 2026.

\bibitem{i128}
Issue 128, ``fix(lower): a method chain on a function call becomes a
constant naming missing variables,'' project repository
\code{HDL/txhdl}, 2026.

\bibitem{i131}
Issue 131, ``feat(test): co-simulate a unit of several clocks, and
lower a clock-crossing channel,'' project repository \code{HDL/txhdl},
2026.

\bibitem{cover}
F.~Filmar and D.~Janković, ``TxHDL documents,'' project repository
\code{HDL/txhdl}, \code{//docs:cover}, 2026.

\end{thebibliography}

\end{document}
